\documentclass[12pt]{article}
\usepackage[T1]{fontenc}
\usepackage[utf8]{inputenc}
\usepackage{lmodern}
\usepackage{textcomp}
\usepackage{enumitem}
\usepackage{geometry}
\geometry{margin=1in}
\begin{document}
\begin{center}
Answer Key - Constitutional Law - Undergraduate Level Assessment 10 \
\small (Answers are ordered exactly as in the question paper.)
\end{center}

\section*{Multiple Choice}
\begin{enumerate}[label=\arabic*.]
\item \textbf{Answer:} B
\\ \textit{Options:} A) The Council of Europe was established in 1950 and consists of 27 member states; B) The Council of Europe was established in 1949 and consists of 47 member states; C) The Council of Europe was established in 1959 and consists of 34 member states; D) The Council of Europe was established in 1984 and consists of 19 member states
\\ \textit{Note:} The correct answer is accurate because the Council of Europe was indeed established in 1949 and currently comprises 47 member states, making it a significant entity in promoting human rights, democracy, and the rule of law in Europe.
\item \textbf{Answer:} D
\\ \textit{Options:} A) The letter is an offer to sell; B) A valid offer cannot be made by letter.; C) The letter contains a valid offer which will terminate within a reasonable time.; D) The letter lacks one of the essential elements of an offer.
\\ \textit{Note:} The letter lacks the essential element of definiteness in an offer because it does not specify the terms of the sale, such as price or quantity, thereby failing to create a binding agreement.
\item \textbf{Answer:} C
\\ \textit{Options:} A) Armed force was prohibited; B) Armed force was permitted with no restrictions; C) Armed force was permitted subject to few restrictions; D) Armed force was not regulated under international law prior to 1945
\\ \textit{Note:} Prior to the United Nations Charter, states operated under the principle of sovereignty that allowed them to use armed force for self-defense or in pursuit of national interests, with limited constraints compared to the post-Charter framework that established stricter regulations on the use of force.
\item \textbf{Answer:} A
\\ \textit{Options:} A) Article 2(4) encompasses only armed force; B) Article 2(4) encompasses all types of force, including sanctions; C) Article 2(4) encompasses all interference in the domestic affairs of States; D) Article 2(4) encompasses force directed only against a State's territorial integrity
\\ \textit{Note:} Article 2(4) of the UN Charter specifically prohibits the use of armed force by member states in their international relations, thereby limiting its scope to military action rather than other forms of force or coercion.
\item \textbf{Answer:} C
\\ \textit{Options:} A) Utilitarianism; B) Communitarianism.; C) Cognitivism.; D) Positivism.
\\ \textit{Note:} Cognitivism supports the idea that ethical statements can be objectively true or false, which challenges ethical relativism's stance that moral truths are culturally bound, thereby providing a strong argument for the universality of human rights.
\end{enumerate}

\section*{True / False}
\begin{enumerate}[label=\arabic*.]
\setcounter{enumi}{5}
\item \textbf{Answer:} True
\\ \textit{Options:} A) True; B) False
\\ \textit{Note:} The statement is true because consequentialism evaluates the morality of punishment based on its outcomes, rather than the nature of the crime itself, suggesting that the appropriateness of punishment should be determined by its future consequences rather than its proportionality to the offense.
\item \textbf{Answer:} False
\\ \textit{Options:} A) True; B) False
\\ \textit{Note:} The UN Vienna Declaration 1993 emphasizes the interdependence and indivisibility of all human rights, affirming that civil and political rights are equally important as third generation rights, rather than prioritizing one over the other.
\item \textbf{Answer:} True
\\ \textit{Options:} A) True; B) False
\\ \textit{Note:} Normative legal theory is concerned with examining the ethical and moral principles that underpin legal systems, thus emphasizing how laws should reflect societal values and political ideals.
\item \textbf{Answer:} True
\\ \textit{Options:} A) True; B) False
\\ \textit{Note:} Radical feminists argue that liberal feminists' pursuit of equality reinforces existing patriarchal structures by encouraging women to adapt to male-defined norms rather than challenging and transforming those norms themselves.
\item \textbf{Answer:} False
\\ \textit{Options:} A) True; B) False
\\ \textit{Note:} Critical Race Theorists do not reject concepts like justice; rather, they seek to redefine and expand these ideas to better address and include the lived experiences of people of color, highlighting systemic inequalities and the need for social change.
\end{enumerate}

\section*{Long Form Response}
\begin{enumerate}[label=\arabic*.]
\setcounter{enumi}{10}
\item \textbf{Answer:} The Kadi judgment significantly impacted the incorporation of UN Security Council (UNSC) resolutions by asserting that such resolutions must be interpreted in alignment with human rights principles. This ruling emphasized that, while the UNSC holds considerable authority in maintaining international peace and security, its actions cannot override fundamental human rights. The European Court of Justice highlighted the necessity of balancing state obligations with individual rights, thus reinforcing the primacy of human rights in legal interpretations. Consequently, the Kadi case set a precedent that mandates a human rights-based approach to the implementation of UNSC measures, ensuring that individual liberties are protected even in the context of international security. This shift indicates a growing recognition of the interplay between international law and human rights within constitutional frameworks.
\\ \textit{Note:} correct because it captures the Kadi judgment's assertion that UN Security Council resolutions must be compatible with fundamental human rights, thereby establishing a critical balance between state authority and individual rights within international law.
\item \textbf{Answer:} Armed violence by non-State actors may constitute an armed attack under Article 51 of the UN Charter when it reaches a certain threshold of severity and is conducted in a manner that poses a significant threat to a State's sovereignty or territorial integrity. The Caroline case illustrates this principle, where the United States justified its military action against British-supported insurgents in Canada by arguing that the threat posed was imminent and required immediate self-defense. In this context, non-State actors can be seen as engaging in an armed attack if their actions are coordinated, organized, and result in substantial harm or destruction. Thus, while Article 51 primarily addresses State-to-State conflict, the Caroline case provides a framework for recognizing non-State actors as potential aggressors under specific conditions.
\\ \textit{Note:} correct because it identifies that armed violence by non-State actors must reach a significant threshold of severity and pose an imminent threat to a state's sovereignty to qualify as an armed attack under Article 51 of the UN Charter, as exemplified by the United States' justification in the Caroline case.
\end{enumerate}

\vfill
\noindent\textit{End of Answer Key}
\end{document}